\documentclass[11pt]{article}

\usepackage{amsmath}
\usepackage{booktabs}
\usepackage{graphicx}
\usepackage[hidelinks]{hyperref}
\usepackage{url}

\title{CPSS: A Severity Scale for Cyber-Physical Incidents\\
in Industrial Control Systems}

\author{
  Shohei Suzuki \\
  Independent Researcher \\
  Tokyo, Japan \\
  \texttt{dx94ke.tmu@gmail.com} \\
}

\begin{document}

\maketitle

\begin{abstract}
Cyberattacks targeting Industrial Control Systems (ICS) can produce
physical accidents, environmental contamination, and large-scale
societal disruption. Existing metrics evaluate either cybersecurity
vulnerabilities (e.g., CVSS) or process safety risks (e.g., HAZOP, SIL),
but no general severity scale exists for cyber incidents that produce
physical consequences.

This paper introduces the Cyber-Physical Severity Scale (CPSS),
a logarithmic scale designed to quantify the impact of cyber incidents
affecting industrial systems. CPSS evaluates four dimensions:
human casualties, environmental damage, infrastructure disruption,
and economic loss. The scale distinguishes between actual and potential
severity, enabling comparison of both realized and worst-credible
consequences across incidents.

We present a dataset of 100 documented ICS cyber incidents scored
under the CPSS framework, spanning the period 2000--2024,
each entry annotated with a primary source reference,
source type (academic, government, vendor report, or news),
and a confidence level reflecting the completeness and
consistency of available public information.
\end{abstract}

%-------------------------------------------------------------------
\section{Introduction}
%-------------------------------------------------------------------

Industrial Control Systems (ICS) operate critical infrastructure
including power plants, water treatment facilities, oil and gas
pipelines, chemical plants, and transportation systems.
Cyberattacks targeting these systems may produce physical consequences:
explosions, toxic releases, equipment destruction, and infrastructure
disruption affecting large populations.

Several incidents demonstrate this risk:
Stuxnet (2010) destroyed uranium enrichment centrifuges at Iran's
Natanz facility~\cite{langner};
the Ukraine power grid attacks (2015, 2016) left hundreds of thousands
of customers without electricity~\cite{eisac};
Triton/TRISIS (2017) disabled safety instrumented systems at a
petrochemical plant in Saudi Arabia~\cite{fireeye};
the Colonial Pipeline ransomware (2021) disrupted fuel supply across
the southeastern United States.

Despite the growing threat, there is no standard severity scale for
cyber incidents that produce physical industrial consequences.
Existing metrics address cybersecurity vulnerabilities or industrial
process hazards independently, but do not capture the full chain from
cyber intrusion to societal impact.
The absence of such a scale makes it difficult to compare incidents,
prioritize defensive investments, and communicate cyber-physical risk
to policymakers and the public.

This work proposes CPSS to fill this gap.

%-------------------------------------------------------------------
\section{Related Work}
%-------------------------------------------------------------------

The Common Vulnerability Scoring System (CVSS)~\cite{cvss} quantifies
the severity of software vulnerabilities based on exploitability and
impact on confidentiality, integrity, and availability.
CVSS does not address physical consequences of cyber incidents.

Process safety methodologies such as Hazard and Operability Study
(HAZOP)~\cite{hazop} and Safety Integrity Levels (SIL/IEC~61511)
analyze industrial process hazards.
These methods address physical risk within normal operational contexts,
but do not incorporate cyber attack vectors.

The International Nuclear Event Scale (INES)~\cite{ines} classifies
nuclear accidents by radiological impact on a seven-level ordinal scale.
INES provides a useful precedent for domain-specific severity
classification, but applies only to nuclear events.

The IEC~62443 standard series~\cite{iec62443} addresses security for
industrial automation and control systems through a zone-and-conduit
model that defines security levels (SL~0--4) based on attack resistance.
IEC~62443 focuses on vulnerability and resilience requirements;
it does not quantify the physical consequence severity of a successful
attack. CPSS is designed to complement IEC~62443 by providing a
consequence metric that can be applied once an attack scenario and
its physical impact are assessed.

CPSS extends the severity-scale approach across all ICS sectors,
integrating the cyber intrusion pathway with physical impact assessment.

More broadly, CPSS belongs to the family of disaster severity scales
used to classify large-scale natural and technological events.
Examples include the Richter magnitude scale for earthquakes,
the Volcanic Explosivity Index (VEI) for volcanic eruptions~\cite{vei},
and the International Nuclear and Radiological Event Scale (INES)
for nuclear accidents.
These scales provide a common language for comparing events across
different locations and time periods, enabling communication of risk
to policymakers, engineers, and the public.
CPSS extends this concept to cyber-physical incidents in
industrial control systems.
As industrial infrastructure becomes increasingly digitized,
cyber-physical incidents represent an emerging class of disasters
that require comparable severity classification.

%-------------------------------------------------------------------
\section{Cyber-Physical Incident Model}
%-------------------------------------------------------------------

ICS cyber incidents follow a causal chain from intrusion to impact:

\begin{equation*}
\text{Cyber Attack} \to
\text{Control Manipulation} \to
\text{Process Deviation} \to
\text{Hazard Release} \to
\text{Impact} \to
\text{CPSS Score}
\end{equation*}

Control manipulation causes process variables (temperature, pressure,
flow, level, speed, or voltage) to deviate from safe operating ranges.
Process deviation in the HAZOP sense~\cite{hazop}---more, less, none,
reverse, or unstable---releases hazardous energy or materials.
The resulting exposure and propagation determine the final impact.

%-------------------------------------------------------------------
\section{Hazard Taxonomy}
%-------------------------------------------------------------------

CPSS classifies physical hazards released by ICS cyber incidents
into six categories, with a seventh value of \emph{none} for
incidents in which no physical process was manipulated:

\begin{itemize}
  \item \textbf{Chemical} --- toxic or flammable substance release
  \item \textbf{Thermal} --- fire, explosion, or heat release
  \item \textbf{Mechanical} --- equipment failure, structural damage
  \item \textbf{Electrical} --- arc flash, grid instability
  \item \textbf{Radiological} --- radioactive material release
  \item \textbf{Hydrological} --- uncontrolled water or sewage release
\end{itemize}

When no physical hazard is triggered (e.g., reconnaissance-only
intrusions), the hazard is recorded as \emph{none}.

%-------------------------------------------------------------------
\section{Impact Dimensions}
%-------------------------------------------------------------------

CPSS evaluates four dimensions of impact, each scored on a 0--10 scale.

\subsection{Human Impact (H)}

\begin{equation}
  H = \log_{10}(100 \cdot \text{deaths} + 10 \cdot \text{hospitalized} + 1)
  \label{eq:H}
\end{equation}

The $\times 100$ multiplier reflects the qualitatively greater severity
of a fatality relative to an equivalent-count infrastructure disruption.
The hospitalized weight ($\times 10$) reflects approximately one-tenth of a
fatality, consistent with disability-adjusted life year (DALY) disability
weights for serious injury requiring hospitalisation~\cite{who_daly}.
$H = 0$ when deaths and hospitalized are both zero; $H > 0$ for any
confirmed fatality or hospitalisation.
Under this formulation, $H = 4$ corresponds to approximately 100 deaths
or 1\,000 hospitalisations.
Table~\ref{tab:H} shows representative reference values.

\begin{table}[htbp]
\centering
\caption{Human Impact scale}
\label{tab:H}
\begin{tabular}{rrl}
\toprule
$H$ & Deaths & Hospitalized (deaths $= 0$) \\
\midrule
0 & 0 & 0 \\
2 & $\sim$1 & $\sim$10 \\
3 & $\sim$10 & $\sim$100 \\
4 & $\sim$100 & $\sim$1{,}000 \\
5 & $\sim$1{,}000 & $\sim$10{,}000 \\
6 & $\sim$10{,}000 & $\sim$100{,}000 \\
7 & $\sim$100{,}000 & $\sim$1{,}000{,}000 \\
\bottomrule
\end{tabular}
\end{table}

\textbf{Calibration rationale.}
The $\times 100$ multiplier is anchored to the value of statistical life (VSL).
Using VSL $\approx \$10$M (per U.S.\ EPA guidance~\cite{epa_vsl}),
100 fatalities imply a societal cost of approximately \$1B,
corresponding to $C = 4$.
This creates cross-dimensional consistency:
$H = 4$ (100~deaths), $I = 4$ (10~million people with service disruption),
and $C = 4$ (\$1B direct loss) are calibrated to represent equivalent severity.
The implicit \$100 per-person disruption cost embedded in the $H/I$ calibration
is a conservative lower bound for a major infrastructure outage
affecting essential services over several days.
In practice, $H$ takes values in $\{0\} \cup [2,\, 10]$:
a single confirmed fatality ($\text{deaths} = 1$) yields
$H = \log_{10}(101) \approx 2.0$, so the interval $(0,\, 2)$ is unreachable
for integer death counts.

\subsection{Environmental Impact (E)}

\begin{equation}
  E = \log_{10}\!\left(\frac{A \cdot P}{10}\right)
  \qquad (E = 0 \text{ if no contamination release})
  \label{eq:E}
\end{equation}

where $A$ is the contaminated area in square metres and $P$ is a persistence
factor representing estimated remediation duration.
$E$ is capped at 10; values yielding a negative logarithm are treated as $E = 0$.

$P$ takes one of four discrete values anchored to established remediation
literature: $P = 1$ (days to weeks, natural recovery, per WHO
drinking-water guidelines~\cite{who_water}); $P = 10$ (months to one year,
per EPA SPCC oil spill response guidance~\cite{epa_spcc}); $P = 100$
(one to ten years, based on the EPA Superfund average cleanup duration
of approximately seven years~\cite{epa_cercla}); or $P = 1000$
(decades or irreversible, per IAEA post-accident
decontamination standards~\cite{iaea_decon}).

When incident-specific remediation data are unavailable, Table~\ref{tab:P_hazard}
provides default $P$ values by hazard category.
Observed remediation data take precedence over defaults.

\begin{table}[htbp]
\centering
\caption{Default persistence factor $P$ by hazard category and
         representative $E$ calibration values}
\label{tab:P_hazard}
\begin{tabular}{llccc}
\toprule
Hazard & Default $P$ & Example $A$ (m$^2$) & $E$ & Scenario \\
\midrule
Mechanical   & 1    & 100       & 1 & Equipment debris, days \\
Electrical   & 1--10 & 100      & 1--2 & Transformer oil spill, weeks \\
Hydrological & 10   & 10{,}000  & 4 & Local waterway, months \\
Thermal      & 10--100 & 1{,}000 & 3--5 & Hydrocarbon soil, months--years \\
Chemical     & 100  & 1{,}000   & 5 & Industrial release, years \\
Radiological & 1000 & 100       & 4 & Facility contamination, decades \\
\bottomrule
\end{tabular}
\end{table}

\subsection{Infrastructure Impact (I)}
\label{sec:I}

\begin{equation}
  I = \log_{10}\!\left(\frac{\text{population directly affected}}{1{,}000} + 1\right)
  \label{eq:I}
\end{equation}

The divisor of 1,000 calibrates I relative to H: at score 4,
approximately 100 deaths (H) corresponds to 10 million people
with service disruption (I), reflecting the greater per-person
severity of fatality versus temporary loss of service.

\noindent\textbf{Note on hazard = none entries ($I_\text{access}$).}
For ICS-scoped reconnaissance incidents where no physical process
was manipulated (\textit{hazard} = \textit{none}), no population
experienced a service disruption, so the standard population-based
formula does not apply. Instead, a substitute indicator
$I_\text{access}$ is recorded in the \textit{I} field:

\begin{equation}
  I_\text{access} = \log_{10}(\text{compromised ICS systems} + 1)
  \label{eq:I_access}
\end{equation}

This measures the breadth of attacker foothold rather than
societal service disruption. The \texttt{i\_basis} field in the
dataset CSV distinguishes the two interpretations
(\texttt{population} vs.\ \texttt{ics\_systems}).
Users should not compare $I$ and $I_\text{access}$ values directly
across entries without accounting for this distinction.
See the dataset scope discussion in Section~\ref{sec:discussion}.

\subsection{Economic Impact (C)}

\begin{equation}
  C = \log_{10}\!\left(\frac{\text{loss (USD)}}{100{,}000}\right)
  \qquad (C = 0 \text{ if loss} = 0)
  \label{eq:C}
\end{equation}

This gives $C = 3$ for \$100M and $C = 4$ for \$1B. Only direct, confirmed
losses are used; indirect or estimated downstream losses are excluded.
Losses below \$100{,}000 yield a negative log value and are treated as $C = 0$.

%-------------------------------------------------------------------
\section{CPSS Score}
%-------------------------------------------------------------------

\begin{equation}
  \text{CPSS} = \max(H,\; E,\; I,\; C)
  \label{eq:cpss}
\end{equation}

Many disaster scales are defined by the dominant physical consequence
of the event rather than the sum of heterogeneous effects.
The Richter scale reflects peak seismic magnitude;
INES is determined by the dominant radiological impact.
CPSS follows the same principle by defining severity as the maximum
impact dimension.
This approach captures the key characteristic of disasters:
overall severity is determined by the worst consequence,
not by the average or sum of heterogeneous impacts.

An important empirical observation from the incident dataset is that
H rarely dominates the CPSS score for ICS cyber incidents.
Of the 100 incidents in the dataset, $H > 0$ in only two cases
(Chatsworth Train Collision, 2008: $H = 3.6$;
Taum Sauk Dam SCADA Overflow, 2005: $H = 1.6$);
in both, another dimension exceeds $H$
($C = 3.7$ and $E = 6.6$, respectively),
so $H$ is the dominant dimension in zero incidents.
This reflects a characteristic of ICS attacks: adversaries generally
seek infrastructure disruption rather than direct casualties,
and physical safety systems often prevent the worst human outcomes
even when process integrity is compromised.
Consequently, I and C tend to be the dominant dimensions.

\subsection{Vector String Notation}

To preserve multi-dimensional information that is otherwise lost
when a scalar score is computed via the maximum operator,
CPSS defines a compact vector string:

\begin{center}
\texttt{CPSS:1.0/H:\textit{h}/E:\textit{e}/I:\textit{i}/C:\textit{c}/A:\textit{a}/P:\textit{p}}
\end{center}

where \textit{h}, \textit{e}, \textit{i}, \textit{c} are the four
dimension scores rounded to one decimal place,
\textit{a} is CPSS-Actual, and \textit{p} is CPSS-Potential.
The prefix \texttt{CPSS:1.0} identifies the specification version.

\noindent\textbf{Example} (Maroochy Shire, 2000~\cite{slay}):
\texttt{CPSS:1.0/H:0.0/E:4.0/I:0.3/C:0.0/A:4.0/P:5.0}

The vector string enables precise incident comparison and
machine-readable data exchange without sacrificing the
interpretability of the scalar class label.

\subsection{Actual and Potential Severity}

CPSS distinguishes two evaluations:

\begin{itemize}
  \item \textbf{CPSS-Actual}: score based on consequences that occurred.
  \item \textbf{CPSS-Potential}: score based on the worst credible
        consequence given the attack scenario and access achieved,
        corresponding to the safety-engineering concept of
        Maximum Credible Accident.
\end{itemize}

The gap between actual and potential severity is analytically
significant: a large gap indicates that defenses, luck, or attacker
restraint prevented the worst outcome.

To reduce subjectivity in CPSS-Potential assessment, analysts should
anchor their evaluation to the highest access level achieved by the
attacker, using Table~\ref{tab:potential_rubric} as a guide.
The table provides a recommended upper bound on CPSS-Potential
for each access level; analysts may assign lower values when
the target system's hazard profile does not support a higher consequence.
Incident-specific evidence (e.g., confirmed SIS bypass, operator
testimony) takes precedence over rubric defaults.

\begin{table}[htbp]
\centering
\caption{CPSS-Potential assessment rubric by attacker access level}
\label{tab:potential_rubric}
\begin{tabular}{lcc}
\toprule
Attacker Access Level & Max Credible & Basis \\
                      & CPSS-Potential & \\
\midrule
IT network only (OT isolated)              & $\leq 4$ & No direct process control \\
OT-adjacent (HMI, historian, DMZ)          & $\leq 5$ & Limited process influence \\
OT control system reached (PLC/DCS)        & $\leq 6$ & Process deviation achievable \\
Safety system reached (SIS/SIL)            & $\leq 7$ & Last-line defense intact \\
Safety system disabled or bypassed         & $\leq 8$ & No automated safeguard \\
Multiple facilities or national grid       & $\leq 9$ & Cascading failure possible \\
\bottomrule
\end{tabular}
\end{table}

\subsection{Severity Classes}

\begin{table}[htbp]
\centering
\caption{CPSS Severity Classes}
\label{tab:class}
\begin{tabular}{cll}
\toprule
Score & Class & Example consequence \\
\midrule
0         & NONE         & No measurable impact \\
$>0$--2   & LOW          & Localised, contained \\
$>2$--4   & MODERATE     & Facility-level disruption \\
$>4$--6   & SEVERE       & Regional infrastructure disruption \\
$>6$--8   & CATASTROPHIC & National-scale impact \\
$>8$--10  & CRITICAL     & Global or irreversible \\
\bottomrule
\end{tabular}
\end{table}

%-------------------------------------------------------------------
\section{Case Studies}
%-------------------------------------------------------------------

Table~\ref{tab:cases} presents eight representative incidents
from the full dataset of 100 events (2000--2024).
Scores for H, E, I, C reflect CPSS-Actual.

\begin{table*}[t]
\centering
\caption{Selected ICS cyber incidents scored under CPSS}
\label{tab:cases}
\resizebox{\textwidth}{!}{
\begin{tabular}{llccccccc}
\toprule
Event & Year & Sector & Hazard & $H$ & $E$ & $I$ & $C$
      & CPSS$_\text{act}$ / CPSS$_\text{pot}$ \\
\midrule
Maroochy Shire Sewage    & 2000 & Water        & Hydrological & 0 & 4 & 0.3 & 0 & 4 / 5 \\
Stuxnet                  & 2010 & Nuclear      & Mechanical   & 0 & 3 & 1.0 & 2 & 3 / 6 \\
Ukraine Grid 1           & 2015 & Power        & Electrical   & 0 & 1 & 2.4 & 3 & 3 / 6 \\
Triton TRISIS            & 2017 & Oil \& Gas   & Thermal      & 0 & 1 & 0.3 & 2 & 2 / 7 \\
Colonial Pipeline        & 2021 & Oil \& Gas   & none         & 0 & 0 & 3.0 & 3 & 3 / 5 \\
Predatory Sparrow Steel  & 2022 & Manufacturing& Thermal      & 0 & 2 & 0.3 & 2 & 2 / 5 \\
Industroyer2             & 2022 & Power        & Electrical   & 0 & 0 & 0.0 & 0 & 0 / 6 \\
Volt Typhoon ICS$^\dagger$ & 2024 & Power      & none         & 0 & 0 & 0.0 & 1 & 1 / 7 \\
\bottomrule
\end{tabular}
}
\end{table*}
\noindent$^\dagger$ Suspected ICS pre-positioning; confirmed OT compromise not publicly disclosed as of the time of writing.

Several observations emerge from the dataset.

\textbf{High potential, low actual.}
Triton TRISIS (CPSS$_\text{act}=2$, CPSS$_\text{pot}=7$) and
Industroyer2 (CPSS$_\text{act}=0$, CPSS$_\text{pot}=6$) represent
incidents where defenses or attacker restraint prevented catastrophic
outcomes. The gap between actual and potential severity is the largest
in the dataset and underscores the unique risk posed by attacks on
safety-critical systems.

\textbf{IT-to-OT cascade with substantial infrastructure impact.}
Colonial Pipeline (CPSS$_\text{act}=3$, CPSS$_\text{pot}=5$) demonstrates
how an IT-layer ransomware attack can produce significant infrastructure
disruption through proactive operator shutdown, without direct OT system
compromise. The remaining gap of 2 reflects that a more sophisticated
attack reaching OT systems could have produced worse consequences.

\textbf{Physical damage without confirmed fatalities.}
Predatory Sparrow (2022)~\cite{claroty_sparrow} is among the few incidents
in the dataset with publicly documented physical equipment damage
(molten steel spill, worker evacuations), yet CPSS$_\text{act} = 2$
due to the localised scale and the absence of confirmed fatalities ($H = 0$).
This illustrates that physical consequences do not necessarily translate
to high human impact scores.

\textbf{High-potential foiled attack.}
Industroyer2 (2022)~\cite{eset_industroyer2} was detected and neutralised
before execution, resulting in CPSS$_\text{act} = 0$ despite the malware
being designed to issue direct commands to high-voltage substations.
The CPSS$_\text{pot} = 6$ reflects that successful execution could have
produced a regional grid disruption comparable to the 2016 Kyiv blackout.

%-------------------------------------------------------------------
\section{Discussion}
\label{sec:discussion}
%-------------------------------------------------------------------

\textbf{Limitations.}
CPSS scores depend on the quality of available incident data.
Many ICS incidents are not publicly disclosed, and reported figures
for casualties, affected populations, and economic losses are often
incomplete or inconsistent across sources.
The CPSS-Potential score requires expert judgment of counterfactual
scenarios, introducing subjectivity.
Table~\ref{tab:potential_rubric} provides an access-level rubric
to anchor assessments, but two analysts evaluating the same incident
may still produce different potential scores,
particularly for incidents where attacker access is partially disclosed.
An event that severely impacts all four dimensions simultaneously
receives the same CPSS score as one that affects a single dimension
at the same level, since CPSS is designed as a worst-case severity
indicator rather than a compound risk index.
Users requiring compound severity analysis may compute an auxiliary indicator
such as the count of dimensions exceeding a chosen threshold.

\textbf{Dataset scope and composition.}
The 42-incident dataset applies a principled inclusion criterion:
an entry is included if the attack was scoped to ICS or OT systems,
regardless of whether physical consequences occurred.
Specifically, the dataset covers three categories:
(1)~incidents with confirmed physical consequences
(process manipulation, equipment damage, or environmental release);
(2)~IT-to-OT cascade incidents, where an IT-layer attack caused
physical shutdown or operational disruption;
and (3)~ICS-scoped reconnaissance and pre-positioning campaigns,
where attackers reached OT-adjacent systems without triggering
a physical process deviation.
Pure IT incidents---ransomware or data-theft campaigns where OT systems
remained isolated and unaffected---are excluded.

This categorisation is reflected in the \textit{hazard} field of the
dataset schema. Entries in category~(3) carry \textit{hazard = none}
and low CPSS-Actual scores, while their CPSS-Potential scores reflect
the physical consequence achievable given the access obtained.
The gap between CPSS-Actual and CPSS-Potential is itself analytically
meaningful for these entries: Triton TRISIS (actual~=~2,
potential~=~7) and INCONTROLLER/PIPEDREAM (actual~=~1, potential~=~7)
demonstrate that sophisticated actors can reach safety-critical systems
without triggering a physical accident, preserving latent risk
for future exploitation.
Users of the dataset should interpret entries with
\textit{hazard = none} with appropriate caution, as CPSS-Actual
reflects access and remediation costs rather than physical impact.

\textbf{Comparison with existing scales.}
CVSS focuses on software vulnerability characteristics prior to
exploitation and does not address physical impact.
INES is restricted to nuclear events and uses a fixed seven-level
ordinal scale. CPSS is designed as a general-purpose scale applicable
across all ICS sectors, using quantitative logarithmic dimensions
(H, I, E, C) derived from measurable physical quantities, enabling
consistent scoring across all industrial hazard types without
requiring domain-specific expert judgment for any single dimension.

\textbf{Future work.}
Because CPSS-Actual is determined by formula, inter-rater reliability
questions concern only CPSS-Potential.
Structured reliability studies comparing CPSS-Potential assessments across
independent analysts using the access-level rubric would establish the
empirical variance of this dimension and identify rubric bounds that require
refinement.
Additional directions include validating default persistence factors
against remediation records from documented industrial accidents,
extending dataset coverage through community contribution,
and integrating CPSS scoring into existing ICS incident reporting
templates to enable longitudinal severity tracking.

\textbf{Determinism of CPSS-Actual.}
CPSS-Actual is fully deterministic: given the four numeric inputs---death toll,
contaminated area and persistence factor, population with service disruption,
and direct economic loss---the score is uniquely determined by formula.
No analyst judgment is required for the actual score.
This eliminates inter-rater variance for confirmed-consequence scoring;
any analyst applying the same confirmed data will produce the same CPSS-Actual
score without coordinated training or calibration exercises.
Subjectivity enters only in CPSS-Potential, which the access-level rubric
in Table~\ref{tab:potential_rubric} is designed to constrain by bounding
assessments to the highest attacker access level confirmed.

\textbf{Adoption pathway.}
CPSS is designed to spread through use rather than through top-down mandate.
Three properties support organic adoption.
First, the determinism property described above means that any analyst scoring
the same incident with the same confirmed data will arrive at the same
CPSS-Actual score, giving the scale the reproducibility needed for
comparative use across reports and organisations.
Second, the specification, dataset, and reference implementation are
released as open resources under an MIT licence.
Any incident report author may compute and publish a CPSS vector string
alongside their narrative description; any tool may integrate the
reference implementation.
The GitHub repository serves as the canonical, version-controlled source
for the specification and dataset schema.
Third, formal adoption by standards bodies
(ISA99/IEC~TC65, ENISA, CISA) remains a long-term option but is not
a prerequisite for practical use.
Precedent exists for widely-used scales---including CVSS itself---that
originated as research proposals and achieved operational adoption through
community uptake before formal standardisation.

%-------------------------------------------------------------------
\section{Conclusion}
%-------------------------------------------------------------------

This paper introduced CPSS, a severity scale for cyber incidents
affecting industrial control systems.
CPSS evaluates four impact dimensions---human, environmental,
infrastructure, and economic---on a common logarithmic scale,
and defines overall severity as the dominant dimension.
The distinction between CPSS-Actual and CPSS-Potential enables
analysis of both realised consequences and worst-credible outcomes.

A dataset of 100 documented incidents spanning 2000--2024 demonstrates
the framework's applicability across sectors including power, water,
oil and gas, manufacturing, and transportation.
The dataset, reference implementation, and full specification are
released as an open resource to support further research in
cyber-physical risk assessment:
\url{https://github.com/fununnn/cpss}

%-------------------------------------------------------------------
\bibliographystyle{plain}
\begin{thebibliography}{9}

\bibitem{iec62443}
ISA/IEC,
\emph{IEC 62443: Security for Industrial Automation and Control Systems},
International Electrotechnical Commission, Geneva, 2018--2023.

\bibitem{cvss}
FIRST.org,
\emph{Common Vulnerability Scoring System v3.1 Specification},
2019.
\url{https://www.first.org/cvss/specification-document}

\bibitem{hazop}
IEC 61882,
\emph{Hazard and Operability Studies (HAZOP Studies) ---
Application Guide},
International Electrotechnical Commission, 2016.

\bibitem{ines}
International Atomic Energy Agency,
\emph{INES: The International Nuclear and Radiological Event Scale
User's Manual}, 2nd~ed.,
IAEA, Vienna, 2013.

\bibitem{langner}
R.~Langner,
``Stuxnet: Dissecting a Cyberwarfare Weapon,''
\emph{IEEE Security \& Privacy}, vol.~9, no.~3, pp.~49--51, 2011.

\bibitem{eisac}
Electricity Information Sharing and Analysis Center (E-ISAC),
\emph{Analysis of the Cyber Attack on the Ukrainian Power Grid},
March 2016.

\bibitem{fireeye}
FireEye,
\emph{TRITON: New ICS Malware Targeting Safety Instrumented Systems},
December 2017.

\bibitem{vei}
C.~G.~Newhall and S.~Self,
``The Volcanic Explosivity Index (VEI): An Estimate of Explosive
Magnitude for Historical Volcanism,''
\emph{Journal of Geophysical Research: Oceans}, vol.~87, no.~C2,
pp.~1231--1238, 1982.

\bibitem{slay}
J.~Slay and M.~Miller,
``Lessons Learned from the Maroochy Water Breach,''
in \emph{Critical Infrastructure Protection}, IFIP vol.~253,
Springer, 2007, pp.~73--82.

\bibitem{epa_vsl}
U.S. Environmental Protection Agency,
\emph{Mortality Risk Valuation},
\url{https://www.epa.gov/environmental-economics/mortality-risk-valuation},
2024.

\bibitem{who_daly}
World Health Organization,
\emph{Global Health Estimates: Life Expectancy and Leading Causes of Death
and Disability},
WHO, Geneva, 2024.

\bibitem{who_water}
World Health Organization,
\emph{Guidelines for Drinking-water Quality}, 4th ed.,
WHO, Geneva, 2017.

\bibitem{epa_spcc}
U.S. Environmental Protection Agency,
\emph{Spill Prevention, Control, and Countermeasure (SPCC) Rule},
40~CFR Part~112, EPA, Washington, DC, 2024.

\bibitem{epa_cercla}
U.S. Environmental Protection Agency,
\emph{Superfund Remedy Report}, 16th ed.,
EPA-542-R-20-001, EPA, Washington, DC, 2020.

\bibitem{iaea_decon}
International Atomic Energy Agency,
\emph{Cleanup of Large Areas Contaminated as a Result of a Nuclear Accident},
IAEA-TECDOC-1162, IAEA, Vienna, 2000.

\bibitem{claroty_sparrow}
Claroty Team82,
``Predatory Sparrow's Attack on Steel Industry,''
Claroty Research Blog, February 2022.
\url{https://claroty.com/team82/research/predatory-sparrow}

\bibitem{eset_industroyer2}
ESET Research,
``Industroyer2: Industroyer reloaded,''
ESET White Paper, April 2022.
\url{https://www.welivesecurity.com/2022/04/12/industroyer2-industroyer-reloaded/}

\end{thebibliography}

\end{document}
